\section{Problem Formulation}
\label{sec:problem}

\subsection{The platform as a black-box channel}

Let $\X$ be a source video. The uploader applies a pre-compensation $\Gt$ with parameters
$\theta\in\Theta$ and uploads $\U=\Gt(\X)$. The platform returns the playback that a viewer receives,
\begin{equation}
  \Y = \fp(\U) = R \circ C_{\tau(\U)} \circ D_{\tau(\U)}(\U),
  \label{eq:platform}
\end{equation}
where $\tau(\U)$ is the \emph{ladder policy} that selects a tier from observable properties of the
upload (resolution, bitrate, possibly content and popularity), $D_\tau$ is the tier's ingest
filtering (e.g.\ denoising), $C_\tau$ is the tier's encoder with its rate control, and $R$ resamples
to the rendition served to the viewer. The uploader can only query $\fp$; $\tau$, $D$ and $C$ are
unknown, $\tau$ is piecewise constant, and $C$ includes quantisation, so $\fp$ is neither known nor
differentiable. Each query costs a real upload.

Playback quality is measured against the source with a full-reference metric $M(\Y,\X)$ (higher is
better), after resampling $\Y$ to the resolution of $\X$. The uploader's problem is
\begin{equation}
  \theta^\star = \argmax_{\theta\in\Theta}\; J_M(\theta),\qquad
  J_M(\theta) = \E_{\X\sim\mathcal{D}}\big[M\big(\fp(\Gt(\X)),\X\big)\big],
  \label{eq:objective}
\end{equation}
over a distribution $\mathcal{D}$ of source videos. In practice the expectation is an average over a
finite set of clips.

\subsection{Tier gain and filter gain}

We split the parameters into a \emph{tier} component $s$ (here the upload scale factor, which
changes $\tau(\U)$) and a \emph{filter} component $\varphi$ (sharpening, grain, colour, denoising),
$\theta=(s,\varphi)$, with identities $s_0$ (no rescaling) and $\varphi_0$ (no filtering).

\begin{definition}[Gain decomposition]
\label{def:decomposition}
For parameters $\theta=(s,\varphi)$, the total gain over a naive upload splits as
\begin{equation}
  \underbrace{J_M(s,\varphi)-J_M(s_0,\varphi_0)}_{\text{total gain}}
  = \underbrace{J_M(s,\varphi_0)-J_M(s_0,\varphi_0)}_{\text{tier gain}}
  + \underbrace{J_M(s,\varphi)-J_M(s,\varphi_0)}_{\text{filter gain at tier } s}.
  \label{eq:decomposition}
\end{equation}
\end{definition}

The split is exact, but it depends on the order: the filter gain at tier $s$ generally differs from
the filter gain at $s_0$, because the filter interacts with the tier's encoder and ingest filtering.
We report the tier-first order, since the tier is discrete and its effect is large, and we measure
the filter gain at the native tier separately (\cref{sec:results-p0}). A pre-compensation is useful
beyond the creator heuristic only if its filter gain is positive.

\subsection{Objectives, judges and Goodhart's law}
\label{sec:goodhart}

The metric $M$ is a model of perceived quality $Q$. Write $M = Q + \varepsilon$, where
$\varepsilon(\theta)$ is the metric's error at the playback produced by $\theta$. The following
statement is a version of the optimiser's curse~\cite{smith2006optimizer}; we state it because it
explains why the gains we report on the optimised metric are biased upwards.

\begin{proposition}[Selection inflates the optimised metric]
\label{prop:curse}
Let $\Theta_n\subset\Theta$ be the finite set of candidates evaluated by a search, let
$\hat\theta=\argmax_{\theta\in\Theta_n} J_M(\theta)$, and let
$\theta^\dagger=\argmax_{\theta\in\Theta_n} J_Q(\theta)$ be the truly best candidate. If the metric
error has zero mean at $\theta^\dagger$, $\E[\varepsilon(\theta^\dagger)]=0$ (for instance, because
the error is independent of $Q$ and unbiased at every fixed $\theta$), then
\begin{equation}
  \E\big[\varepsilon(\hat\theta)\big] \;\ge\; \E\big[J_Q(\theta^\dagger)-J_Q(\hat\theta)\big] \;\ge\; 0 .
\end{equation}
\end{proposition}

\begin{proof}
By the definition of $\hat\theta$, $J_Q(\hat\theta)+\varepsilon(\hat\theta) \ge
J_Q(\theta^\dagger)+\varepsilon(\theta^\dagger)$, hence $\varepsilon(\hat\theta) \ge
\varepsilon(\theta^\dagger) + J_Q(\theta^\dagger)-J_Q(\hat\theta)$. Take expectations; the last
difference is non-negative by the definition of $\theta^\dagger$.
\end{proof}

The proposition covers errors that are noise-like. A more harmful case is an error with structure:
if $\varepsilon$ grows along some direction of $\Theta$ (for example, a metric that rewards a
contrast increase), the optimiser moves along that direction and the inflation is no longer bounded
by the noise level. This is the regime documented for VMAF and VMAF
NEG~\cite{zvezdakova2019hacking,siniukov2021hacking}.

We therefore distinguish the \emph{objective} $M$ from a set of \emph{judges}
$\mathcal{J}=\{M_1,\dots,M_K\}$, metrics that the optimiser never sees.

\begin{definition}[Judge-confirmed gain]
\label{def:confirmed}
A gain of configuration $a$ over configuration $b$ is \emph{judge-confirmed} if at least $k$ of the
$K$ judges assign $a$ a higher mean score than $b$ over the evaluation clips. In P0 we use $K=3$ and
$k=2$; in P1 we additionally require a confidence interval over clips that excludes zero
(\cref{sec:decision}).
\end{definition}

\begin{remark}
If the judges' errors were independent of each other and of the objective's error, the probability
that a gain which is spurious for $Q$ appears on $k$ of $K$ judges would shrink geometrically in
$k$. Real metrics share failure modes (several react to sharpness or contrast in the same way), so
agreement among judges is evidence, not proof; a viewer study remains the final arbiter.
\end{remark}

\subsection{Optimising through a learned emulator}
\label{sec:emulator-theory}

For a learned pre-compensation, $\Gt$ is a network with millions of parameters, too many for
black-box search. We fit an emulator $\Em\approx\fp$ on pairs $\{(\U_i,\fp(\U_i))\}$ collected from the
black box, and train $\Gt$ by gradient ascent on the emulated objective
\begin{equation}
  J_M^{E}(\theta)=\E_{\X}\big[M\big(\Em(\Gt(\X)),\X\big)\big].
\end{equation}
If $M$ is $L$-Lipschitz in its first argument (with respect to a norm $\|\cdot\|$), then
\begin{equation}
  J_M(\theta) \;\ge\; J_M^{E}(\theta) - L\,\E_{\X}\big\|\fp(\Gt(\X))-\Em(\Gt(\X))\big\|.
  \label{eq:emulator-bound}
\end{equation}
The emulator error in the last term is small only on the distribution of uploads that $\Em$ was
trained on. Gradient ascent on $J_M^E$ moves $\Gt(\X)$ away from that distribution, towards inputs
where the emulator overestimates $M$, so the gap can grow without limit. This is the failure mode of
offline model-based optimisation~\cite{trabucco2021conservative,yu2020mopo}. We measure the gap
directly, by sending $\Gt(\X)$ through the real black box and reporting the emulator's PSNR against
the real playback, and we control it with a trust region:
\begin{equation}
  \min_\theta\; -J_M^{E}(\theta) + \lambda\,\E_{\X}\|\Gt(\X)-\X\|_1
  \quad\text{s.t.}\quad \|\Gt(\X)-\X\|_\infty\le\delta,
  \label{eq:trust}
\end{equation}
which keeps the uploads near the identity-like uploads in the emulator's training data.

A further option, not used in the experiments reported here, is a straight-through
estimator~\cite{bengio2013estimating,athalye2018obfuscated,pesnel2026scaled}: use the real playback
in the forward pass and the emulator only for the gradient,
\begin{equation}
  \tilde\Y = \fp(\U) + \Em(\U) - \sg\big(\Em(\U)\big),
\end{equation}
where $\sg$ is the stop-gradient operator. The objective's value is then always the real one, and
only its gradient is approximate.

\subsection{Content-scoped claims and gating}
\label{sec:content-theory}

The platform's damage depends systematically on content: rate control and ingest filtering respond
to texture, motion and flat regions. A pre-compensation that is not useful on average may still be
useful on a subset. Let $c(\X)$ be a class label computed from content features (\cref{sec:setup}),
with classes $k=1,\dots,K_c$ and class probabilities $p_k$, and let the class-conditional gain of
configuration $a$ over a baseline $b$ be
\begin{equation}
  \Delta_k(a,b)=\E\big[M(\fp(a(\X)),\X)-M(\fp(b(\X)),\X)\,\big|\,c(\X)=k\big].
\end{equation}
A \emph{gate} $\pi$ assigns one configuration to each class. Its gain over the baseline is
$\sum_k p_k\,\Delta_k(\pi(k),b)$, which is at least the gain of the best single configuration if
$\pi(k)$ is chosen per class with the true $\Delta_k$.

The difficulty is estimation. If the classes, the gate or the claim are chosen on the same clips
that are used to estimate $\Delta_k$, the estimate is inflated by selection, exactly as in
\cref{prop:curse}, with the choice of subgroup playing the role of the search. We therefore:
\begin{enumerate}
  \item define the classes from content features by a fixed rule, before any learned model is
        trained;
  \item split the clips into train, development and confirmation sets, grouped by source so that
        no source contributes to two splits;
  \item choose the gate and the claimed class on the development set only;
  \item evaluate the confirmation set once.
\end{enumerate}
This is a pre-registered design in the sense of~\cite{nosek2018preregistration}. The classes,
splits and decision rules are committed to the repository before training (\cref{sec:setup}).

\subsection{Decision rule}
\label{sec:decision}

The unit of analysis is the clip: seeds are averaged within a clip, and uncertainty is estimated
over clips with a percentile bootstrap~\cite{efron1993introduction} (10{,}000 resamples). For the
content-scoped hypothesis, a learned pre-compensation is said to help on class $k$ if, on the
confirmation clips of that class:
\begin{enumerate}
  \item at least two judges show a mean gain over the class-tuned sharpening baseline whose 95\%
        bootstrap interval excludes zero (97.5\% if a second class is tested);
  \item the mean SSIMULACRA2 gain is non-negative;
  \item the sign of the effect holds for at least two of three training seeds.
\end{enumerate}
A judge that is part of the network's training loss is not used to judge it.
